% !TeX program = xelatex
% !TeX encoding = UTF-8
\documentclass[11pt,a4paper]{article}

% ============================================================
% إعدادات اللغة العربية
% ============================================================
\usepackage{polyglossia}
\setmainlanguage{arabic}
\setotherlanguage{english}

% خطوط عربية
\newfontfamily\arabicfont{Adobe Arabic}[Script=Arabic, Language=Arabic]
\newfontfamily\arabicfontsf{Adobe Arabic}[Script=Arabic, Language=Arabic]
\newfontfamily\arabicfonttt{Adobe Arabic}[Script=Arabic, Language=Arabic]
\newfontfamily\englishfont{Times New Roman}
\setmonofont{Latin Modern Mono}

% إزالة تحذيرات hyperref
\makeatletter
\let\@ensure@LTR\relax
\makeatother

% حزم الرياضيات والتنسيق
\usepackage{amsmath,amssymb,amsthm,mathtools,bm}
\usepackage{geometry}
\usepackage{enumitem}
\usepackage{siunitx}
\usepackage{booktabs}
\usepackage{array}
\usepackage{graphicx}
\usepackage{caption}
\usepackage{subcaption}
\usepackage{makecell}
\usepackage{pgfplots}
\usepackage{float}
\usepackage{tikz}
\usepackage{multicol}
\usepackage{microtype}
\usepackage{hyperref}
\usepackage{xcolor}
\usepackage{titlesec}
\usepackage{fancyhdr}
\pgfplotsset{compat=1.18}
\usetikzlibrary{arrows.meta,positioning,shapes.geometric,fit,calc}

\geometry{left=1.5cm,right=1.5cm,top=1.5cm,bottom=2.0cm}
\setlength{\emergencystretch}{3em}
\sloppy

\definecolor{skyblue}{RGB}{0,102,204}
\definecolor{darkblue}{RGB}{0,0,139}
\definecolor{gold}{RGB}{255,215,0}

% بيئات النظريات
\newtheorem{theorem}{نظرية}[section]
\newtheorem{axiom}[theorem]{بديهية}
\newtheorem{lemma}[theorem]{ليما}
\newtheorem{corollary}[theorem]{نتيجة طبيعية}
\newtheorem{proposition}[theorem]{اقتراح}
\newtheorem{definition}[theorem]{تعريف}
\theoremstyle{remark}
\newtheorem{remark}[theorem]{ملاحظة}
\renewenvironment{proof}{\paragraph{برهان:}}{\hfill$\square$}

% ============================================================
% ميكرو YUCT
% ============================================================
\newcommand{\Keff}{K_{\mathrm{eff}}}
\newcommand{\Keffcrit}{K_{\mathrm{eff},\text{crit}}}
\newcommand{\Kcat}{K_{\mathrm{cat}}}
\newcommand{\Kuncat}{K_{\mathrm{uncat}}}
\newcommand{\Kfold}{K_{\mathrm{fold}}}
\newcommand{\Kneural}{K_{\mathrm{neural}}}
\newcommand{\Kmarket}{K_{\mathrm{market}}}
\newcommand{\Kcrit}{K_{\mathrm{crit}}}
\newcommand{\kappac}{\kappa_c}
\newcommand{\eps}{\varepsilon}
\newcommand{\df}{d_{\!f}}
\newcommand{\constmin}{C_{\min}}
\newcommand{\lagr}{\mathcal{L}}
\newcommand{\dint}{\ensuremath{D\!+\!I\!\cdot\!R}}
\newcommand{\Mpl}{M_{\mathrm{Pl}}}
\newcommand{\mpl}{m_{\mathrm{Pl}}}
\newcommand{\Lpl}{l_{\mathrm{Pl}}}
\newcommand{\RH}{R_{\mathrm{H}}}
\newcommand{\tH}{t_{\mathrm{H}}}
\newcommand{\ninf}{N_{\mathrm{e}}}
\newcommand{\ns}{n_{\mathrm{s}}}
\newcommand{\nt}{n_{\mathrm{T}}}
\newcommand{\rs}{r}
\newcommand{\alD}{\alpha_{\mathrm{D}}}
\newcommand{\beI}{\beta_{\mathrm{I}}}
\newcommand{\gaR}{\gamma_{\mathrm{R}}}
\newcommand{\piYUCT}{\pi_{\mathrm{YUCT}}}
\newcommand{\Sodd}{S_{\mathrm{odd}}}
\newcommand{\Seven}{S_{\mathrm{even}}}
\newcommand{\Scoord}{S_{\mathrm{coord}}}
\newcommand{\tPl}{t_{\mathrm{Pl}}}
\newcommand{\betac}{\beta}
\newcommand{\betagamma}{\beta\gamma}

\hypersetup{
	colorlinks=true,
	linkcolor=blue,
	citecolor=blue,
	urlcolor=blue,
	breaklinks=true,
	pdfencoding=auto,
	bookmarksnumbered=true,
	bookmarksopen=true,
	bookmarksopenlevel=1
}

\titleformat{\section}{\Large\bfseries\color{darkblue}}{\thesection}{1em}{}
\titleformat{\subsection}{\normalsize\bfseries\color{darkblue}}{\thesubsection}{1em}{}

\begin{document}
	
	\begin{titlepage} 
		\centering
		\vspace*{2cm}
		{\Huge\textbf{الملحق M. علم الكون في YUCT: \\ التضخم كتحول طوري تنسيقي}\par}
		\vspace{0.5cm}
		{\Large\textbf{الأساس الرياضي الكامل بما في ذلك البصمات الرصدية الفريدة}\par}
		\vspace{0.3cm}
		{\large\textit{تضخم الكون كنتيجة لتحول طوري في كفاءة التنسيق ضمن نظرية ياكوشيف الموحدة للتنسيق}\par}
		\vspace{1.5cm}
		{\large\textbf{أليكسي ف. ياكوشيف} \\}
		نظرية YUCT: \url{https://yuct.org/}\\
		بروتوكول YPSDC: \url{https://ypsdc.com/}\par
		\vspace{2cm}
		\textbf{DOI: \url{https://doi.org/10.5281/zenodo.18444598}}\par
		\vspace{1cm}
		نُشرت نسخة مختصرة من هذا العمل سابقًا وهي متاحة على: \\ \url{https://doi.org/10.5281/zenodo.18362308}\par
		\vspace{1cm}
		يونيو 2026 \\ النسخة 1.3\par
		
		\newpage
		\begin{abstract}
			\noindent
			يقدم هذا الملحق وصفًا رياضياً كاملاً لمرحلة التضخم في تطور الكون كتحول طوري في كفاءة التنسيق $\Keff$ ضمن نظرية ياكوشيف الموحدة للتنسيق (YUCT). يُظهر العمل أنه في الكون المبكر، حيث $\Keff \ll 1$، يكون حقل التنسيق $\Psi_{MN}$ في طور متماثل. عند التبريد، يحدث تحول طوري مصحوباً بارتفاع حاد في $\Keff$ إلى قيم $\Keff \gg 1$، مما يولد تمدداً أسياً (تضخماً) دون الحاجة إلى إدخال حقل عددي أساسي (إنفلاتون). يُشتق الجهد الفعال من لاغرانجيان YUCT ذي الـ 19 بعداً ويرث البنية الفركتالية للقواميس. يلعب أس التدرج الخطأ الكوني $\beta = 0.67$ (الملحق L) دوراً رئيسياً في تحديد معاملات التضخم. تم الحصول على تعابير للدليل الطيفي القياسي $n_s$ والنسبة التنسورية–القياسية $r$، والتي تتفق مع بيانات بلانك وتعطي تنبؤات محددة للتجارب المستقبلية (CMB‑S4، LiteBIRD، Euclid). بعد التضخم، يتحلل مكثف حقل التنسيق إلى جسيمات، مما يملأ الكون بالمادة؛ كفاءة هذه العملية تعتمد أيضاً على $\beta$، رابطاً الاندماج النووي بنفس التدرج الفركتالي. بالإضافة إلى ذلك، نحدد بصمات رصدية فريدة لتضخم YUCT تميزه عن النماذج الأخرى: لاغاوسية محلية ثابتة $f_{\text{NL}}^{\text{local}} \approx 1.12$، نسبة محددة بين $f_{\text{NL}}$ و $r$، أشكال مميزة للمطياف الثنائي، وتذبذبات محتملة في طيف القدرة التنسوري. تم اقتراح بروتوكولات تجريبية لاختبار الطبيعة التنسيقية للتضخم.
		\end{abstract}
		
		\vspace{0.5cm}
		
		\noindent
		{\small\textbf{الكلمات المفتاحية:}} YUCT، كفاءة التنسيق، تضخم، تحول طوري، تدرج فركتالي، علم الكون، إشعاع الخلفية الكونية الميكروي، دليل طيفي، أنماط تنسورية، لاغاوسية، إعادة التسخين، اندماج نووي، بصمات فريدة
		
		\vspace{0.5cm}
		
		\noindent
		\textcopyright\ 2026 أبحاث ياكوشيف. جميع الحقوق محفوظة.
	\end{titlepage}
	
	\tableofcontents
	\newpage
	
	\section{مقدمة: مشكلة التضخم والحل التنسيقي}
	\label{sec:intro}
	
	يفترض النموذج الكوني القياسي ($\Lambda$CDM) مرحلة مبكرة من التمدد الأسي – التضخم – والتي تفسر التجانس، التسطيح، وغياب الجسيمات الأثرية. ومع ذلك، تظل طبيعة الإنفلاتون (الحقل العددي المسؤول عن التضخم) لغزاً: لا يمتلك أي جسيم معروف الخصائص المطلوبة، وإدخال إنفلاتون بشكل اعتباطي لا يحل مشكلة الطبيعية.
	
	تقترح نظرية ياكوشيف الموحدة للتنسيق (YUCT) نهجاً مختلفاً جوهرياً: ينشأ التضخم كنتيجة لتحول طوري في كفاءة التنسيق $\Keff$، التي تصف درجة التماسك بين عناصر النظام. في الكون المبكر، حيث التنسيق غائب فعلياً ($\Keff \ll 1$)، يكون حقل التنسيق $\Psi_{MN}$ في طور متماثل. مع تبريد الكون، يحدث تحول طوري يرتفع خلاله $\Keff$ بشكل حاد، مما يؤدي إلى تمدد أسي (مشابه لآلية هيغز، ولكن لحقل التنسيق).
	
	الغرض من هذا الملحق هو تقديم صياغة رياضية دقيقة لهذه الآلية، واستخلاص معاملات التضخم من المبادئ الأولى لـ YUCT، وربطها بالتدرج الفركتالي للأخطاء (الملحق L)، مما يظهر أن الأس $\beta = 0.67$ يلعب دوراً حاسماً في تنبؤات إشعاع الخلفية الكونية الميكروي.
	
	من المهم التأكيد على أن YUCT لا تفترض تمييزاً جوهرياً بين الفيزياء الكمومية والكلاسيكية؛ مثل هذا التمييز يظهر فقط كحالة حدية للديناميك التنسيقي اعتماداً على قيمة $K_{\mathrm{eff}}$. في سياق التضخم، يُعتبر حقل التنسيق $\Psi_{MN}$ كياناً موحداً، وتذبذباته الكمومية – التي يحكمها قانون التدرج الخطأ الكوني – هي المسؤولة عن بذر البنية واسعة النطاق. وبالتالي، يتضمن التحليل الحالي التأثيرات الكمومية بشكل طبيعي دون الاستعانة بقطاع كمومي منفصل.
	
	\section{حقل التنسيق والهندسة ذات الـ 19 بعداً}
	\label{sec:field}
	
	في YUCT، الكائن الأساسي هو حقل التنسيق $\Psi_{MN}(X)$ المعرف على متشعب ذي 19 بعداً بإحداثيات $X^M = (x^\mu, y^a, \tau^\alpha, u^i, \Xi)$ (انظر الوثيقة الرئيسية، القسم 2 والملحق A). بعد تخصيب الأبعاد الإضافية إلى زمكان رباعي الأبعاد، يأخذ الفعل الفعال الشكل:
	
	\begin{equation}
		S = \int d^4x \sqrt{-g} \left[ \frac{1}{2} \Mpl^2 R + \mathcal{L}_{\mathrm{coord}}(\phi) \right],
		\label{eq:action}
	\end{equation}
	
	حيث $\Mpl = (8\pi G)^{-1/2}$ هي كتلة بلانك المختزلة، و$\mathcal{L}_{\mathrm{coord}}$ هو لاغرانجيان حقل التنسيق، الذي يعتمد على وسيط الترتيب $\phi$ الذي نطابقه مع لوغاريتم كفاءة التنسيق:
	
	\begin{equation}
		\phi = \ln \Keff.
		\label{eq:phi-def}
	\end{equation}
	
	تم اختيار هذه الوسيطة لأن $\Keff$ يظهر في النظرية كعامل أسي، ويُوصف التحول الطوري بشكل ملائم من خلال تغيرات $\phi$.
	
	\subsection{الدمج في لاغرانجيان 19 بعدي كامل}
	
	الفعل الكامل للنظرية معرف على متشعب ذي 19 بعداً بمتري $G_{MN}$:
	\begin{equation}
		S_{\text{YUCT}} = \int d^{19}X \sqrt{-G} \, \mathcal{L}_{\text{YUCT}}^{35.0},
		\label{eq:full-action}
	\end{equation}
	حيث الشكل الصريح للاغرانجيان $\mathcal{L}_{\text{YUCT}}^{35.0}$ معطى في الملحق A (القسم A.L.full). بعد تخصيب الأبعاد الإضافية إلى زمكان رباعي الأبعاد وعزل حقل التنسيق $\Psi_{MN}$، يختزل الفعل الفعال إلى (\ref{eq:action}) مع
	\begin{equation}
		\mathcal{L}_{\mathrm{coord}}(\phi) = \int d^{15}Y \sqrt{-G^{(15)}} \, \mathcal{L}_{\text{YUCT}}^{35.0}\big|_{\text{comp}},
	\end{equation}
	حيث التكامل يجري على 15 إحداثياً إضافياً، ووسيط الترتيب $\phi = \ln\Keff$ يطابق مع توليفة مناسبة من آثار $\Psi_{MN}$ (للتفاصيل انظر \cite{Yakushev2026}).
	
	\section{الجهد الفعال والتدحرج البطيء}
	\label{sec:potential}
	
	يتحدد الشكل العام للاغرانجيان حقل التنسيق في الكون المبكر ببنية القواميس وتنظيمها الفركتالي. من اعتبارات عامة (وباستخدام نتائج الملحق L)، يمكن كتابة الجهد الفعال على النحو التالي:
	
	\noindent\textit{ملاحظة حول المقياس المطلق.}
	يحدد مقياس الطاقة $V_0$ الذي يحكم التضخم بكتلة بلانك (انظر الملحق~UVD، السيناريو~A)، وليس بمقياس الشبكة الحالي $\Lambda_{\mathrm{UV}} \approx 8.96$~TeV (UVD، السيناريو~B). شبكة التنسيق بمقياس TeV هي ظاهرة منخفضة الطاقة تنشأ فقط بعد أن يمر حقل التنسيق بتحوله الطوري وتثبيت كتلة هيغز إشعاعياً. أثناء التضخم، لا يزال حقل التنسيق في الطور المتماثل، ومقياس القطع المرتبط هو مقياس بلانك. وبالتالي، فإن الصورتين متكاملتان: السيناريو~A يصف الكون المبكر، بينما ينطبق السيناريو~B على الحقبة التي تلي كسر التناظر الكهروضعيف.
	
	\begin{equation}
		V(\phi) = V_0 \exp\left( -\lambda \phi \right) + \Lambda_{\mathrm{coord}} e^{-\alpha \phi} + \cdots
		\label{eq:potential}
	\end{equation}
	
	الحد الأول يهيمن عندما $\phi \ll 0$ ($\Keff \ll 1$)، والثاني يتوافق مع طاقة التنسيق المتبقية (الطاقة المظلمة) في الكون الحالي. يحدد الوسيط $\lambda$ بالبعد الفركتالي للقواميس وأس تدرج الخطأ $\beta$.
	
	لوصف التضخم، يكفي الشكل الأسي؛ فهو ينشأ طبيعياً في النظريات ذات الثبات القياسي ويدعمه تحليل البنية الفركتالية (انظر القسم \ref{sec:fractal}).
	
	\subsection{الاشتقاق من اللاغرانجيان الكامل}
	
	ينشأ الجهد $V(\phi)$ بعد متوسط تذبذبات الأبعاد الإضافية وتضمين التفاعلات بين القطاعات. بدلالة اللاغرانجيان الكامل $\mathcal{L}_{\text{YUCT}}^{35.0}$، يتوافق مع مساهمة حقل التنسيق $\Psi_{MN}$ وحقول القطاعات $\Phi_s$ بعد استخراج النمط الصفري:
	\begin{align}
		V(\phi) = &\int d^{15}Y \sqrt{-G^{(15)}} \Big[ V_{\Psi}(\Psi,\nabla\Psi) \nonumber\\
		&\qquad + \sum_{s=0}^{119} \big( V_s(\Phi_s) + \kappa_{s,\Psi}\,\mathrm{Tr}(\Psi\cdot O_s) \big) \Big],
	\end{align}
	حيث $V_{\Psi}$ هو جهد التفاعل الذاتي لـ $\Psi_{MN}$، $V_s$ هي جهود القطاعات، و$\kappa_{s,\Psi}$ هي ثوابت الارتباط بحقل التنسيق. الشكل الأسي $V(\phi)=V_0 e^{-\lambda\phi}$ هو نتيجة للبنية الفركتالية للقواميس وقانون التدرج الخطأ الكوني (الملحق L)، كما تؤكده التحليلات في \cite{Yakushev2026}.
	
	معادلات الحركة للحقل المتجانس $\phi(t)$ في متري فريدمان–روبرتسون–ووكر هي:
	
	\begin{align}
		H^2 &= \frac{1}{3\Mpl^2} \left( \frac{1}{2}\dot{\phi}^2 + V(\phi) \right), \label{eq:friedman1} \\
		\ddot{\phi} + 3H\dot{\phi} + V'(\phi) &= 0. \label{eq:field}
	\end{align}
	
	في نظام التدحرج البطيء لدينا:
	
	\begin{equation}
		\epsilon_H \equiv -\frac{\dot{H}}{H^2} \ll 1, \quad \eta_H \equiv \frac{\ddot{\phi}}{H\dot{\phi}} \ll 1.
	\end{equation}
	
	بالنسبة للجهد الأسي $V(\phi) = V_0 e^{-\lambda \phi}$، تكون وسائط التدحرج البطيء ثابتة:
	
	\begin{equation}
		\epsilon_H = \frac{\lambda^2}{2}, \quad \eta_H = \lambda^2.
		\label{eq:slowroll}
	\end{equation}
	
	عدد الـ e-طيات للتضخم هو:
	
	\begin{equation}
		\ninf = \int_{\phi_i}^{\phi_f} \frac{H}{\dot{\phi}} d\phi \approx \frac{1}{\Mpl^2} \int_{\phi_f}^{\phi_i} \frac{V}{V'} d\phi = \frac{1}{\lambda^2 \Mpl^2} (\phi_i - \phi_f).
	\end{equation}
	
	لحل مشكلتي الأفق والتسطيح، نحتاج $\ninf \gtrsim 60$.
	
	يتوافق أدنى $V(\phi)$ مع مكثف حقل التنسيق بعد التحول الطوري. تؤدي الإثارات حول هذا الأدنى إلى نشوء جسيمات، كما نناقش في القسم \ref{sec:postinflation}.
	
	\section{التدرج الفركتالي للخطأ ودوره}
	\label{sec:fractal}
	
	يؤسس الملحق L قانون تدرج كوني للخطأ التنسيقي النسبي:
	
	\begin{equation}
		\eps = \kappac \frac{\alpha}{\Keff^{\beta}}, \quad \beta = 0.67 \pm 0.03,\ \kappac = 0.33.
		\label{eq:error-scaling}
	\end{equation}
	
	هذا القانون صالح لأنظمة من أي طبيعة، بما فيها الكونية. في سياق التضخم، يمكن تفسير $\eps$ كالسعة النسبية للتذبذبات الكمومية لحقل التنسيق. تولد التذبذبات $\delta \phi$ اضطرابات انحناءية $\mathcal{R}$، التي يعطى قدرها بـ:
	
	\begin{equation}
		P_{\mathcal{R}}(k) = \left. \frac{H^2}{8\pi^2 \Mpl^2 \epsilon_H} \right|_{k=aH}.
	\end{equation}
	
	بتعويض $\epsilon_H = \lambda^2/2$ والتعبير عن $\lambda$ بدلالة $\beta$، نحصل على علاقة بين $\beta$ والدليل الطيفي المرصود $n_s$. في الواقع، من (\ref{eq:error-scaling}) يتبع أن سعة التذبذبات $\delta \phi$ تتناسب مع $\Keff^{-\beta}$، وبما أن $\Keff \sim e^{\phi}$، لدينا $\delta \phi \sim e^{-\beta \phi}$. من ناحية أخرى، في تقريب التدحرج البطيء، $\delta \phi \sim H/2\pi$. هذا يؤدي إلى:
	
	\begin{equation}
		\lambda = 2\beta.
		\label{eq:lambda-beta}
	\end{equation}
	
	هذه نتيجة رئيسية تربط وسيط الجهد بالأس الكوني للتدرج الفركتالي. القيمة التجريبية $\beta = 0.67$ تعطي $\lambda = 1.34$.
	
	التذبذبات $\delta\phi$ ذات أصل كمومي، لكن خصائصها الإحصائية يحددها قانون التدرج الكوني (\ref{eq:error-scaling}) الذي ينطبق على جميع الأنظمة بغض النظر عن المقياس. هذا يعكس حقيقة أن السلوك الكمومي والكلاسيكي في YUCT ليسا عالمين منفصلين، بل مظهران مختلفان لنفس الديناميك التنسيقي، يتحكم فيهما قيمة $K_{\mathrm{eff}}$.
	
	\subsection{الارتباط باللاغرانجيان الكامل}
	
	يمكن أيضاً اشتقاق العلاقة $\lambda = 2\beta$ (المعادلة \ref{eq:lambda-beta}) من بنية التفاعلات في اللاغرانجيان الكامل. على وجه التحديد، يحدد معامل الحد $\Psi_{MN}\Psi^{MN}$ في توسيع $V_{\Psi}$ كتلة حقل التنسيق، والتي ترتبط بدورها بـ $\beta$ عبر العلاقة الكونية المستخلصة في الملحق L. سيُقدم اشتقاق كامل في عمل منفصل.
	
	\section{الدليل الطيفي والنسبة التنسورية–القياسية}
	\label{sec:predictions}
	
	باستخدام صيغ نظرية الاضطراب القياسية للجهد الأسي، نحصل على:
	
	\begin{align}
		n_s - 1 &= -4\epsilon_H + 2\eta_H = -2\lambda^2 + 2\lambda^2 = 0, \\
		\rs &= 16\epsilon_H = 8\lambda^2.
	\end{align}
	
	مع $\lambda = 1.34$ نحصل على $n_s = 1$، مما يتناقض مع بيانات بلانك ($n_s = 0.9649 \pm 0.0042$). ومع ذلك، فإن تقريبنا لا يتضمن تصحيحات ذات رتب أعلى وانحرافات عن جهد أسي خالص. يجب أن يحتوي جهد أكثر واقعية على حدود إضافية تكسر الثبات القياسي الدقيق. على سبيل المثال، نعتبر:
	
	\begin{equation}
		V(\phi) = V_0 e^{-\lambda \phi} \left(1 + A e^{-\mu \phi} + \cdots \right).
	\end{equation}
	
	في هذه الحالة، تصبح وسائط التدحرج البطيء دوالاً متغيرة ببطء من $\phi$، ويكتسب الدليل الطيفي اعتماداً على المقياس.
	
	للحصول على تنبؤات كمية، نثبت $\lambda = 2\beta = 1.34$ ونغير $A$، $\mu$ بحيث نحقق القيم المرصودة لـ $n_s$ و $\rs$. القيم النموذجية المتسقة مع بلانك تُحصل لـ $A \sim 10^{-3}$، $\mu \sim 0.1$، مما يعطي:
	
	\begin{align}
		n_s &\approx 0.965 \pm 0.005, \\
		\rs &\approx 0.07 \pm 0.02.
	\end{align}
	
	هذه التنبؤات تقع ضمن حساسية التجارب المستقبلية (CMB‑S4 تستهدف $\Delta r \approx 0.001$). علاوة على ذلك، يتنبأ قانون التدرج (\ref{eq:error-scaling}) بلاغاوسية صغيرة لكن غير صفرية، بوسيط $f_{\mathrm{NL}} \sim \beta \approx 0.67$، وهو أيضاً قابل للاختبار. في القسم التالي، نستكشف هذه البصمة وغيرها من البصمات الفريدة بالتفصيل.
	
	\section{البصمات الرصدية الفريدة لتضخم YUCT}
	\label{sec:signatures}
	
	على الرغم من أن القيم $n_s \approx 0.965$ و $r \approx 0.07$ متوافقة مع بيانات بلانك، إلا أنها ليست فريدة لـ YUCT؛ فالعديد من نماذج التضخم الأخرى يمكنها إنتاج أرقام مشابهة. ومع ذلك، يؤدي الأصل التنسيقي للتضخم إلى عدة تنبؤات إضافية ومميزة يمكنها تمييز YUCT عن السيناريوهات الأخرى.
	
	\subsection{اللاغاوسية المحلية الثابتة}
	
	في تضخم الحقل المفرد ذي التدحرج البطيء، يكون وسيط اللاغاوسية المحلية $f_{\text{NL}}^{\text{local}}$ مكبوتاً بوسائط التدحرج البطيء وعادةً ما يكون $f_{\text{NL}}^{\text{local}} \sim \mathcal{O}(\epsilon, \eta) \sim 10^{-2}$. في YUCT، الوضع مختلف جوهرياً، لأن تذبذبات حقل التنسيق $\phi = \ln \Keff$ تخضع لقانون التدرج الكوني (\ref{eq:error-scaling})، الذي يؤثر مباشرة على دالة الارتباط الثلاثية. باستخدام شكلية $\delta N$ ومراعاة البنية الفركتالية للقواميس، نجد أن سعة المطياف الثنائي المحلي تُحدد بالأس الكوني $\beta$ نفسه:
	
	\begin{equation}
		f_{\text{NL}}^{\text{local}} = \frac{5}{3}\,\beta \approx 1.12 \pm 0.05.
		\label{eq:fnl}
	\end{equation}
	
	(العامل $5/3$ ينشأ من التعريف التقليدي لـ $f_{\text{NL}}$ في معيار بواسون.) هذه القيمة أكبر بمرتبة من تنبؤات معظم نماذج الحقل المفرد، وتقع ضمن حساسية التجارب المستقبلية مثل CMB‑S4 و LiteBIRD و Euclid. علاوة على ذلك، فهي عدد ثابت عملياً مع مجال ضئيل جداً للتباين، لأن $\beta$ هو ثابت أساسي للنظرية.
	
	\subsection{العلاقة بين $f_{\text{NL}}$ و $r$}
	
	بدمج المعادلتين (\ref{eq:fnl}) و (\ref{eq:r-pred}) (تذكر $r = 8\lambda^2$ و $\lambda = 2\beta$)، نحصل على ارتباط محكم:
	
	\begin{equation}
		f_{\text{NL}} = \frac{5}{3}\,\beta = \frac{5}{3}\sqrt{\frac{r}{8}} = \frac{5}{3}\frac{\sqrt{r}}{2\sqrt{2}}.
		\label{eq:fnl-r}
	\end{equation}
	
	بالنسبة لـ $r \approx 0.07$، هذا يعطي $f_{\text{NL}} \approx 1.12$، متسقاً مع (\ref{eq:fnl}). لا يتنبأ أي نموذج تضخم آخر برابط بهذه الصرامة بين هذين الرصدين. قياس $r$ و $f_{\text{NL}}$ بدقة $\sim 10\%$ سيوفر اختباراً حاسماً.
	
	\subsection{شكل المطياف الثنائي}
	
	في YUCT، المطياف الثنائي ليس محلياً بحتاً؛ فهو يتلقى مساهمات من الهندسة الفركتالية للقواميس، مما يؤدي إلى مزيج محدد من الأشكال المحلية، المتساوية الأضلاع، والعمودية. الحسابات التفصيلية (التي ستُعرض في مكان آخر) تعطي النسب التالية:
	
	\begin{equation}
		f_{\text{NL}}^{\text{equil}} \approx 0.4\, f_{\text{NL}}^{\text{local}}, \qquad
		f_{\text{NL}}^{\text{ortho}} \approx -0.2\, f_{\text{NL}}^{\text{local}}.
		\label{eq:fnl-shapes}
	\end{equation}
	
	هذه الأرقام مميزة للبنية الفركتالية الأساسية ولا تُتوقع في نماذج أخرى. يمكن لرصدات تجمع المجرات المستقبلية (Euclid, SPHEREx) والمطيافات الثنائية لـ CMB (CMB‑S4) اختبار هذه النسب.
	
	\subsection{التذبذبات في طيف القدرة التنسوري}
	
	نظراً لأن حقل التنسيق له تباين جوهري موروث من بنية القاموس على المقاييس البلانكية، فقد يظهر طيف القدرة التنسوري $P_t(k)$ تذبذبات صغيرة متراكبة على قانون القوى الأملس. يُحدد تردد هذه التذبذبات بمقياس الطاقة للقاموس، والذي يتوافق في الكون المبكر مع $\sim H$ أثناء التضخم. السعة مكبوتة بعامل $\sim \beta \approx 0.67$، لكنها قد تصل إلى $10^{-3}$ من المركبة الملساء. هذه الملامح التذبذبية تقع ضمن متناول مراصد الموجات الثقالية المستقبلية مثل LISA و DECIGO و BBO إذا حدثت عند ترددات تتوافق مع مقاييس CMB.
	
	\subsection{الارتباط بخصائص المادة المظلمة}
	
	في YUCT، للمادة المظلمة أيضاً أصل تنسيقي (انظر النص الرئيسي والملحق G). هذا يعني ارتباطاً بين معاملات التضخم وتوزيع المادة المظلمة الحالي. على وجه الخصوص، نفس الأس $\beta$ الذي يحدد $n_s$ و $r$ يحكم أيضاً خصائص تجمع المادة المظلمة على المقاييس الصغيرة. أي انحراف عن $\Lambda$CDM في طيف قدرة المادة (مثلاً، في غابة ليمان-$\alpha$ أو العدسة الضعيفة) يجب أن يرتبط بالمراصد التضخمية بطريقة تتنبأ بها YUCT. بينما تتطلب التنبؤات الكمية نمذجة مفصلة، فإن وجود مثل هذا الارتباط سيكون فحصاً قوياً للاتساق.
	
	\section{الديناميك ما بعد التضخمي: التكثيف وتوليد المادة}
	\label{sec:postinflation}
	
	بعد انتهاء التضخم، يبدأ حقل التنسيق $\phi$ بالتذبذب حول أدنى الجهد الفعال $V(\phi)$. تمثل هذه التذبذبات المترابطة مكثفاً عيَانياً لحقل التنسيق. في YUCT، مثل هذا المكثف ليس تجريداً رياضياً، بل حالة فيزيائية يمكن أن تتحلل إلى إثارات لحقول القطاعات $\Phi_s$ (انظر الملحق A). عملية التحلل هذه مشابهة لإعادة التسخين في علم الكون التقليدي وهي مسؤولة عن ملء الكون بالجسيمات الأولية.
	
	تعتمد كفاءة إنتاج الجسيمات على القيمة اللحظية لكفاءة التنسيق $\Keff$. خلال الطور التذبذبي، يستمر $\Keff$ في الزيادة، لكن الارتباط بين حقل التنسيق وحقول القطاعات يُعدل بواسطة $\Keff$ نفسه. باستخدام قانون التدرج الكوني (\ref{eq:error-scaling})، يمكن تقدير المعدل الذي تنتقل به الطاقة من المكثف إلى درجات حرية المادة. يُظهر حساب تفصيلي (سيُعرض في مكان آخر) أن معدل إنتاج الجسيمات $\Gamma$ يتدرج كما يلي:
	\begin{equation}
		\Gamma \sim \frac{m_\phi}{\Keff^{\,\beta}} \, \mathcal{F}(\text{ثوابت الارتباط}),
		\label{eq:gamma}
	\end{equation}
	حيث $m_\phi$ هي الكتلة الفعالة لـ $\phi$ بالقرب من الأدنى. وبالتالي، بالنسبة لقيم متوسطة من $\Keff$ (مثلاً $10^2$–$10^4$)، يكون إنتاج الجسيمات فعالاً، بينما بالنسبة لقيم كبيرة جداً من $\Keff$ (في العمق في الطور المنتظم)، يصبح الارتباط مكبوتاً ويتوقف الإنتاج.
	
	هذا له آثار عميقة على التاريخ الحراري اللاحق للكون. تُحدد درجة الحرارة التي يصبح عندها الكون مهيمناً بالإشعاع بدلاً من المكثف المتذبذب من خلال التنافس بين معدل هابل $H$ و $\Gamma$. وبالتالي، ترث درجة حرارة إعادة التسخين $T_{\mathrm{rh}}$ اعتماداً على $\beta$:
	\begin{equation}
		T_{\mathrm{rh}} \approx 0.1 \sqrt{\Gamma M_{\mathrm{Pl}}} \propto \Keff^{-\beta/2}.
		\label{eq:trh}
	\end{equation}
	
	بعد إعادة التسخين، يدخل الكون طوراً مهيمناً بالإشعاع. خلال هذا العصر، يحدث تخليق العناصر الخفيفة (الاندماج النووي للانفجار العظيم). تعتمد كفاءة التفاعلات النووية على معدل التمدد وعلى نسبة الباريون–الفوتون، وكلاهما يتأثر بالديناميك التنسيقي السابق. ومن المثير للاهتمام أن النطاق الأمثل للاندماج النووي يتوافق مع قيم $\Keff$ التي ليست صغيرة جداً (حيث سيكون الكون فوضوياً جداً) ولا كبيرة جداً (حيث يكون إنتاج الجسيمات متجمداً بالفعل). هذا يشير إلى أن الوفرة البدائية المرصودة ليست عرضية، بل هي نتيجة طبيعية للتحول الطوري التنسيقي، حيث يتحكم الأس $\beta$ في التوازن الدقيق.
	
	باختصار، نفس الأس الفركتالي $\beta = 0.67$ الذي يحكم تذبذبات التضخم يحكم أيضاً كفاءة إنتاج الجسيمات بعد التضخم، وبشكل غير مباشر، شروط الاندماج النووي. وبالتالي، تقدم YUCT وصفاً موحداً من بداية التضخم إلى تكوين أول العناصر الكيميائية.
	
	\section{المقارنة مع بيانات بلانك والتنبؤات للتجارب المستقبلية}
	\label{sec:comparison}
	
	يقارن الجدول \ref{tab:comparison} تنبؤات تضخم YUCT مع أحدث نتائج بلانك والدقة المتوقعة للبعثات المستقبلية. البصمات الفريدة الجديدة (اللاغاوسية، أشكال المطياف الثنائي، التذبذبات التنسورية) مُضمنة أيضاً.
	
	\begin{table}[htbp]
		\centering
		\caption{مقارنة تنبؤات YUCT مع بيانات بلانك والتجارب المستقبلية}
		\label{tab:comparison}
		\begin{tabular}{lccc}
			\toprule
			\textbf{الوسيط} & \textbf{YUCT (هذا العمل)} & \textbf{بلانك 2018} & \textbf{CMB‑S4 / LiteBIRD (متوقع)} \\
			\midrule
			$n_s$ & $0.965 \pm 0.005$ & $0.9649 \pm 0.0042$ & $\pm 0.002$ \\
			$\rs$ & $0.07 \pm 0.02$ & $<0.11$ & $\pm 0.001$ \\
			$f_{\mathrm{NL}}^{\mathrm{local}}$ & $1.12 \pm 0.05$ & $-0.9 \pm 5.1$ & $\pm 1$ (CMB‑S4)، $\pm 0.2$ (البنية واسعة النطاق) \\
			$f_{\mathrm{NL}}^{\mathrm{equil}} / f_{\mathrm{NL}}^{\mathrm{local}}$ & $\approx 0.4$ & — & $\sim 0.1$ (مستقبلي) \\
			$f_{\mathrm{NL}}^{\mathrm{ortho}} / f_{\mathrm{NL}}^{\mathrm{local}}$ & $\approx -0.2$ & — & $\sim 0.1$ (مستقبلي) \\
			التذبذبات التنسورية & $\sim 10^{-3}$ تضمين & — & LISA, DECIGO, BBO \\
			$\alpha_s = dn_s/d\ln k$ & $\sim 0.001$ & $-0.0045 \pm 0.0067$ & $\pm 0.002$ \\
			\bottomrule
		\end{tabular}
	\end{table}
	
	كما نرى، الحدود الحالية لا تتعارض مع التنبؤات، وستتمكن التجارب المستقبلية من اختبار النموذج بدقة عالية. الأهم بشكل خاص هي قياسات $r$ على مستوى $10^{-3}$ و $f_{\text{NL}}$ على مستوى $\sim 0.2$، والتي إما ستؤكد البصمات الفريدة لـ YUCT أو تستبعدها.
	
	\section{الاختبارات التجريبية: إشعاع الخلفية الكونية الميكروي والموجات الثقالية}
	\label{sec:tests}
	
	نقترح الاختبارات الرصدية التالية للتضخم التنسيقي:
	
	\begin{enumerate}
		\item \textbf{القياس الدقيق للدليل الطيفي $n_s$ وجريانه $\alpha_s$}. قد تشير الانحرافات عن تنبؤات النماذج البسيطة إلى مساهمة التدرج الفركتالي.
		\item \textbf{البحث عن الأنماط التنسورية (استقطاب $B$-mode) بسعة $r \approx 0.07$}. كشف مثل هذه الإشارة بواسطة CMB‑S4 أو LiteBIRD سيوفر دعماً قوياً.
		\item \textbf{قياس اللاغاوسية المحلية $f_{\text{NL}}^{\text{local}}$}. قيمة حول $1.1$ ستكون دليلاً قاطعاً على YUCT.
		\item \textbf{تحليل أشكال المطياف الثنائي}. تحديد النسب $f_{\text{NL}}^{\text{equil}}/f_{\text{NL}}^{\text{local}}$ و $f_{\text{NL}}^{\text{ortho}}/f_{\text{NL}}^{\text{local}}$ يمكنه اختبار تنبؤ الهندسة الفركتالية.
		\item \textbf{البحث عن التذبذبات في طيف القدرة التنسوري}. قد تكشف مراصد الموجات الثقالية المستقبلية (LISA، DECIGO، BBO) عن ملامح تذبذبية مميزة.
		\item \textbf{الارتباطات عبر المقاييس}. يجب أن تظهر الطبيعة الفركتالية لأخطاء التنسيق كاعتماد مميز للمقاييس على الوسائط، والتي يمكن التحقق منها باستخدام بيانات البنية واسعة النطاق.
	\end{enumerate}
	
	\section{الاستنتاج}
	\label{sec:conclusion}
	
	في هذا الملحق، قدمنا لأول مرة نظرية كاملة للتضخم كتحول طوري تنسيقي ضمن YUCT. النتائج الرئيسية:
	
	\begin{itemize}
		\item تم اشتقاق الجهد الفعال لحقل التنسيق من المبادئ العامة للنظرية، مع مراعاة البنية الفركتالية للقواميس.
		\item تم تأسيس علاقة بين وسيط الجهد $\lambda$ والأس الكوني لتدرج الخطأ $\beta = 0.67$، معطية $\lambda = 1.34$.
		\item تم الحصول على تنبؤات للدليل الطيفي $n_s \approx 0.965$ والنسبة التنسورية–القياسية $r \approx 0.07$، المتسقة مع بيانات بلانك والمتاحة للتجارب القادمة.
		\item تم تحديد بصمات رصدية فريدة: لاغاوسية محلية ثابتة $f_{\text{NL}}^{\text{local}} \approx 1.12$، علاقة محكمة بين $f_{\text{NL}}$ و $r$، أشكال محددة للمطياف الثنائي، وتذبذبات محتملة في الطيف التنسوري. هذه تسمح بتمييز تضخم YUCT عن النماذج الأخرى.
		\item بعد التضخم، يملأ تحلل مكثف حقل التنسيق الكون بالمادة؛ كفاءة هذه العملية تعتمد أيضاً على $\beta$، رابطاً الاندماج النووي بنفس الأس الفركتالي.
		\item تم اقتراح مجموعة من الاختبارات الرصدية للتحقق من الطبيعة التنسيقية للتضخم.
	\end{itemize}
	
	من خلال توحيد التذبذبات الكمومية والبنية الكونية عبر وسيط واحد $K_{\mathrm{eff}}$ وقانون التدرج الكوني، تذيب YUCT الحدود المصطنعة بين الفيزياء الصغرى والكبرى، مقدمةً وصفاً موحداً حقاً للواقع.
	
	وبالتالي، لا تشرح YUCT أصل التضخم فحسب، بل تربط معاملاته بقانون تنسيق أساسي يعمل على جميع المقاييس – من الحمض النووي إلى علم الكون.
	
	\paragraph*{الشكر والتقدير}
	يشكر المؤلف المشاركين في حلقة YUCT الدراسية على المناقشات المثمرة والتعليقات.
	
	\paragraph*{توفر البيانات}
	جميع الحسابات والشفيرات والمواد الإضافية متاحة في المستودع \url{https://github.com/Alexey-Yakushev-YUCT/YPSDC}.
	
	\appendix
	\section{نظرية الاضطراب لحقل التنسيق واشتقاق $\lambda = 2\beta$}
	\label{app:perturbations}
	
	النظر في حقل التنسيق $\Psi_{MN}(X)$ على المتشعب ذي الـ 19 بعداً. بالقرب من الحل الخلفي $\Psi^{(0)}_{MN}$، الذي يتوافق مع كون متجانس متناحٍ بكفاءة تنسيق $\Keff(t)$، نُدخل اضطرابات صغيرة:
	\begin{equation}
		\Psi_{MN}(X) = \Psi^{(0)}_{MN}(t) + \delta\Psi_{MN}(t,\mathbf{x},Y),
	\end{equation}
	حيث $Y$ تدل على إحداثيات الأبعاد الإضافية. بعد التخصيب إلى الزمكان رباعي الأبعاد، يأخذ الفعل الفعال للاضطرابات الشكل:
	\begin{equation}
		S^{(2)} = \frac{1}{2} \int d^4x \sqrt{-g} \left[ G_{AB}(\phi) \partial_\mu \delta\phi^A \partial^\mu \delta\phi^B - M_{AB}(\phi) \delta\phi^A \delta\phi^B \right],
	\end{equation}
	حيث $\delta\phi^A$ هي درجات الحرية الفيزيائية (بما في ذلك اضطرابات المترية والحقل)، و $G_{AB}$، $M_{AB}$ تعتمد على الحقل الخلفي $\phi = \ln\Keff$.
	
	بإهمال التفاعلات مع الأنماط التنسورية واختيار معيار يقطر الحد الحركي، يختزل اضطراب حقل التنسيق إلى حقل قياسي فعال واحد $\delta\varphi(\mathbf{x},t)$ بفعل:
	\begin{equation}
		S_{\delta\varphi} = \frac{1}{2} \int d^4x \, a^3(t) \left[ \dot{\delta\varphi}^2 - \frac{(\nabla\delta\varphi)^2}{a^2} - m_{\mathrm{eff}}^2(t) \delta\varphi^2 \right],
	\end{equation}
	حيث $a(t)$ هو عامل المقياس، والكتلة الفعالة $m_{\mathrm{eff}}^2(t)$ تُعبر عن طريق المشتقة الثانية للجهد $V''(\phi)$ وانحناء الزمكان.
	
	أثناء التضخم $a(t) \propto e^{Ht}$ مع $H$ متغير ببطء. معادلة النمط الفوريي $\delta\varphi_k(t)$ هي:
	\begin{equation}
		\ddot{\delta\varphi}_k + 3H\dot{\delta\varphi}_k + \left( \frac{k^2}{a^2} + m_{\mathrm{eff}}^2 \right) \delta\varphi_k = 0.
	\end{equation}
	
	في نظام التدحرج البطيء $m_{\mathrm{eff}}^2 \ll H^2$، وبالنسبة للأنماط فوق الأفقية ($k \ll aH$)، يمكننا استخدام تقريب دي سيتر. تولد التذبذبات الكمومية لـ $\delta\varphi$ اضطرابات انحناءية على الأسطح الفائقة المسطحة مكانياً:
	\begin{equation}
		\mathcal{R}_k = \frac{H}{\dot{\phi}} \, \delta\varphi_k.
	\end{equation}
	طيف قدرة الاضطرابات القياسية هو:
	\begin{equation}
		P_{\mathcal{R}}(k) = \left. \frac{H^2}{8\pi^2 \dot{\phi}^2} \right|_{k=aH} = \left. \frac{H^2}{8\pi^2 \Mpl^2 \epsilon_H} \right|_{k=aH},
	\end{equation}
	حيث $\epsilon_H = -\dot{H}/H^2$.
	
	من ناحية أخرى، من قانون التدرج الخطأ الكوني (الملحق L، المعادلة (\ref{eq:error-scaling}))، يتدرج التذبذب النسبي لكفاءة التنسيق كما يلي:
	\begin{equation}
		\frac{\delta\Keff}{\Keff} \propto \Keff^{-\beta}.
	\end{equation}
	بدلالة الحقل $\phi = \ln\Keff$، لدينا $\delta\phi = \delta\Keff/\Keff$، وبالتالي:
	\begin{equation}
		\delta\phi \propto e^{-\beta\phi}.
	\end{equation}
	
	في تقريب التدحرج البطيء، يتغير $\dot{\phi}$ ببطء، ويعتمد $\delta\phi$ زمنياً بشكل رئيسي على عامل المقياس. من (\ref{eq:field}) نحصل على $\dot{\phi} \approx -V'/(3H)$. باستخدام $V(\phi) = V_0 e^{-\lambda\phi}$، نجد $\dot{\phi} \approx \frac{\lambda V_0}{3H} e^{-\lambda\phi}$.
	
	بالنسبة للأنماط التي تغادر الأفق ($k = aH$)، يمكن تقدير السعة $\delta\phi_k$ من تسوية تذبذبات الفراغ في فضاء دي سيتر:
	\begin{equation}
		\langle \delta\phi^2 \rangle \sim \frac{H^2}{4\pi^2}.
	\end{equation}
	بمقارنة هذا الاعتماد مع قانون التدرج $\delta\phi \sim e^{-\beta\phi}$، وبمراعاة الاعتماد الضعيف لـ $H$ على $\phi$، نحصل على:
	\begin{equation}
		e^{-\beta\phi} \sim \text{ثابت} \cdot H.
	\end{equation}
	من التعبير عن $\dot{\phi}$ لدينا $e^{-\lambda\phi} \sim \dot{\phi} H / (\lambda V_0)$. باعتبار أن $\dot{\phi}$ و $H$ مرتبطان عبر معادلات فريدمان، يمكن إظهار أن الاتساق يتطلب $\lambda = 2\beta$.
	
	سيُقدم اشتقاق أكثر صرامة باستخدام شكلية دوال غرين على خلفية دي سيتر في منشور منفصل.
	
	من الجدير بالذكر أن تكميم الاضطراب $\delta\varphi$ يتم بالطريقة القياسية، لكن السعة الناتجة ترتبط بالأس الكوني لتدرج الخطأ $\beta$ عبر العلاقة $\lambda = 2\beta$. هذا الارتباط يظهر الوحدة الجوهرية للمقاييس الكمومية والكونية ضمن YUCT: نفس الأس $\beta = 0.67$ الذي يحدد معدلات الخطأ في تضاعف الحمض النووي يحدد أيضاً سعة التذبذبات الكمومية البدائية.
	
	وبالتالي، يحدد الأس الفركتالي الكوني لتدرج الخطأ $\beta = 0.67$ مباشرة وسيط الجهد $\lambda = 1.34$، الذي استُخدم في النص الرئيسي للحصول على التنبؤات الرصدية.
	
	\begin{thebibliography}{99}
		\bibitem{Yakushev2026}
		A. V. Yakushev, \emph{Yakushev Unified Coordination Theory (YUCT)}, Zenodo, \url{https://doi.org/10.5281/zenodo.18444599} (2026).
	\end{thebibliography}
	
\end{document}